\chapter[Lapisan Implementasi \& Migrasi: Peta Jalan Transformasi]{Lapisan
Implementasi \& Migrasi: Work Package, Plateau, dan Peta Jalan}
\label{ch:implementation-migration}
\setlength{\emergencystretch}{2em}

\section*{Tujuan Pembelajaran}
\begin{itemize}
  \item Menjelaskan elemen lapisan implementasi \& migrasi ArchiMate---\emph{work
        package}, \emph{deliverable}, \emph{implementation event}, \emph{plateau},
        dan \emph{gap}.
  \item Menyusun \textbf{peta jalan} (\emph{roadmap}) transformasi: rangkaian
        plateau dari \emph{baseline} ke \emph{target} beserta paket kerja yang
        mewujudkannya.
  \item Menempatkan analisis dalam TOGAF ADM \textbf{Fase E (Peluang \& Solusi)}
        dan \textbf{Fase F (Perencanaan Migrasi)}.
  \item Menutup rantai keterlacakan seluruh buku: dari motivasi sampai paket kerja
        yang membangun arsitektur, dan mewujudkannya kembali menjadi kapabilitas.
\end{itemize}

Tujuan pertama dan kedua mendukung CLO 3; tujuan ketiga mendukung CLO 2 dan
CLO 6; tujuan keempat mendukung CLO 7 karena mahasiswa merangkai seluruh lapisan
menjadi satu argumen transformasi yang utuh.

\section{Ringkasan Awal Bab}

Enam bab sebelumnya menjawab \emph{apa} arsitektur target: motivasi, strategi,
bisnis, data, aplikasi, dan teknologi. Bab ini menjawab \emph{bagaimana
mencapainya}: bagaimana organisasi berpindah dari keadaan sekarang ke keadaan
target, langkah demi langkah, tanpa harus mengganti semuanya sekaligus. Dalam
ArchiMate, lapisan implementasi \& migrasi memuat \textbf{work package} (paket
kerja/proyek), \textbf{deliverable} (hasil kerja), \textbf{implementation event}
(peristiwa implementasi), \textbf{plateau} (keadaan arsitektur yang stabil pada
suatu periode), dan \textbf{gap} (kesenjangan antar-plateau).

Pola intinya: \textbf{work package mewujudkan deliverable}; \textbf{deliverable
mewujudkan plateau}; \textbf{plateau mencakup} kapabilitas dan requirement yang
aktif pada keadaan itu; dan \textbf{gap} menandai selisih antara \emph{baseline}
dan \emph{target}. Karena ini lapisan terakhir, keluarannya---peta jalan dan
analisis kesenjangan---menjadi bagian \textbf{Rekomendasi} dan \textbf{Rencana}
pada naskah, sekaligus menutup keseluruhan model.

\section{Pendahuluan}

Arsitektur target yang indah tidak berguna jika tidak ada jalan untuk mencapainya.
Kebanyakan kegagalan transformasi bukan karena visinya salah, melainkan karena
lompatannya terlalu besar: anggaran habis, sistem lama tak sempat dibungkus, dan
kepatuhan tertinggal. Lapisan implementasi \& migrasi memaksa kita berpikir
\emph{bertahap}---memecah perjalanan menjadi beberapa \textbf{plateau} yang
masing-masing stabil, bernilai, dan dapat dievaluasi.

Bagi makalah, lapisan ini mengubah rekomendasi arsitektur menjadi \textbf{rencana
yang dapat dinilai kelayakannya}. Pendekatan bertahap dan pengiriman inkremental
sejalan dengan praktik rekayasa perangkat lunak modern \cite{newman2021building},
sedangkan penahapan dan tata kelola arsitektur mengikuti TOGAF ADM Fase E dan F
\cite{togaf10}. Fondasi eksekusi---urutan investasi yang membangun kapabilitas
secara kumulatif---selaras dengan gagasan \emph{operating model} dan
\emph{foundation for execution} \cite{ross2006enterprise}.

\section{Skenario Pembuka dan Pertanyaan Pemandu}

Arsitektur target rel transfer Bank Wanua sudah lengkap sampai lapisan teknologi.
Kini komite investasi bertanya: mulai dari mana? Dengan anggaran terbatas dan
sistem inti lama yang tak bisa diganti sekaligus, tim harus memecah program
menjadi beberapa rilis. Tenggat kepatuhan SNAP BI menambah tekanan: sebagian
kapabilitas harus hidup lebih dulu.

\subsection*{Pertanyaan Pemandu}
\begin{enumerate}
  \item Apa keadaan \emph{baseline} saat ini dan apa \emph{target} akhirnya?
  \item Kesenjangan (gap) apa yang memisahkan keduanya?
  \item Paket kerja apa yang menutup tiap kesenjangan, dan apa deliverable-nya?
  \item Plateau antara apa yang masuk akal sebagai tonggak bernilai?
  \item Bagaimana urutan, ketergantungan, dan peristiwa penting (mis. go-live)
        disusun menjadi peta jalan?
\end{enumerate}

\section{Mengapa Bab Ini Penting bagi Pembaca Non-TI}

Justru pembaca non-TI---manajer, pemilik anggaran, regulator internal---yang
paling sering membaca lapisan ini, karena di sinilah arsitektur bertemu waktu,
biaya, dan risiko. Dengan lima notasi sederhana, seorang analis dapat menjelaskan
\emph{apa} yang dikerjakan (work package), \emph{apa hasilnya} (deliverable),
\emph{kapan keadaan stabil tercapai} (plateau), dan \emph{apa yang masih kurang}
(gap)---tanpa jargon proyek. Peta jalan menjadi bahasa bersama antara arsitek dan
pengambil keputusan.

\section{Konsep Inti dan Glosarium Mini}

\subsection{Lapisan Implementasi \& Migrasi ArchiMate}

\textbf{Work package} adalah rangkaian tindakan dengan awal dan akhir yang jelas
(sebuah proyek atau fase proyek) yang mewujudkan sekumpulan deliverable.
\textbf{Deliverable} adalah hasil kerja yang terdefinisi jelas dari sebuah work
package. \textbf{Implementation event} adalah peristiwa terkait implementasi,
misalnya \emph{kickoff} atau \emph{go-live}. \textbf{Plateau} adalah keadaan
arsitektur yang relatif stabil pada suatu periode. \textbf{Gap} adalah hasil
analisis kesenjangan antara dua plateau---apa yang berbeda antara \emph{baseline}
dan \emph{target}.

\subsection{Notasi Elemen di Archi}

Gambar~\ref{fig:implementation-legend} menampilkan jenis elemen lapisan ini di
Archi. Semua elemen berwarna merah muda (\emph{peach}); ikon di pojok kanan atas
membedakan jenisnya.

\begin{figure}[htbp]
\centering
\includegraphics[width=0.86\textwidth]{implementation-legend.pdf}
\caption{Jenis elemen lapisan implementasi \& migrasi ArchiMate dan simbolnya di
Archi.}
\label{fig:implementation-legend}
\end{figure}

\begin{table}[htbp]
\centering
\scriptsize
\caption{Elemen lapisan implementasi \& migrasi, kegunaan, dan contohnya.}
\label{tab:impl-elements}
\begin{tabularx}{\textwidth}{@{}p{0.22\textwidth} p{0.30\textwidth} X@{}}
\hline
\textbf{Elemen} & \textbf{Kegunaan} & \textbf{Contoh} \\
\hline
Work Package & Proyek/fase dengan awal-akhir jelas yang mewujudkan deliverable. &
WP1 Bangun Platform API \& Gateway; WP2 Bangun Mesin AML. \\
Deliverable & Hasil kerja terdefinisi dari work package. & Platform API \& Gateway
Terpasang; Mesin AML \& Limit Teruji. \\
Implementation Event & Peristiwa implementasi/migrasi. & Kickoff Program; Go-Live
SNAP BI. \\
Plateau & Keadaan arsitektur yang stabil pada suatu periode. & Baseline; MVP Rel
Transfer; Target (SNAP BI Penuh). \\
Gap & Selisih (hasil analisis kesenjangan) antara dua plateau. & Belum ada
gateway; AML masih manual. \\
\hline
\end{tabularx}
\end{table}

\subsection{Plateau, Gap, Work Package, dan Deliverable}

Empat elemen ini bekerja bersama membentuk pola migrasi
(Gambar~\ref{fig:migration-pattern}). Sebuah \textbf{plateau} menandai keadaan
stabil; dua plateau dihubungkan oleh relasi \emph{triggering} (memicu) yang
menyatakan urutan waktu. Selisih di antara keduanya adalah \textbf{gap}. Untuk
menutup gap, sebuah \textbf{work package} dijalankan; ia \textbf{mewujudkan}
(\emph{realization}) sebuah \textbf{deliverable}, dan deliverable itu pada
gilirannya \textbf{mewujudkan} plateau target. Dengan kata lain: gap menyatakan
\emph{apa} yang kurang, work package menyatakan \emph{siapa mengerjakan apa}, dan
plateau menyatakan \emph{keadaan bernilai} yang dicapai.

\begin{figure}[htbp]
\centering
\scalebox{0.95}{
\begin{tikzpicture}[
  font=\footnotesize,
  impl/.style={rectangle, rounded corners=1.5mm, draw=red!55!black, fill=red!8,
    thick, align=center, text width=30mm, minimum height=10mm},
  plat/.style={rectangle, rounded corners=1.5mm, draw=red!55!black, fill=red!14,
    thick, align=center, text width=26mm, minimum height=10mm},
  real/.style={-{Triangle[open]}, thick, draw=gray!70, densely dashed},
  trig/.style={-Latex, thick, draw=gray!70},
  assoc/.style={thick, draw=gray!60}
]
  \node[plat] (p0) {Plateau\\\emph{Baseline}};
  \node[plat, right=42mm of p0] (p2) {Plateau\\\emph{Target}};
  \node[impl, below=16mm of p0] (wp) {Work Package};
  \node[impl, right=14mm of wp] (del) {Deliverable};
  \node[draw=red!55!black, fill=red!8, thick, circle, minimum size=7mm,
    inner sep=0pt, font=\scriptsize] (gap) at ($(p0)!0.5!(p2)$) {};
  \node[font=\scriptsize] at ($(gap)+(0,4mm)$) {Gap};
  \draw[trig] (p0) -- node[above,font=\scriptsize]{memicu} (p2);
  \draw[assoc] (gap) -- ($(gap)+(0,-10mm)$) node[below,font=\scriptsize]{pada};
  \draw[real] (wp) -- node[below,font=\scriptsize]{mewujudkan} (del);
  \draw[real] (del.north) to[out=90,in=-90] node[right,font=\scriptsize]{mewujudkan} (p2.south);
\end{tikzpicture}
}
\caption{Pola migrasi (kanonik ArchiMate): plateau \emph{baseline} memicu plateau
\emph{target}; gap menandai selisihnya; work package \emph{mewujudkan}
deliverable, dan deliverable \emph{mewujudkan} plateau target.}
\label{fig:migration-pattern}
\end{figure}

\subsection{TOGAF ADM Fase E \& F}

Fase E (Peluang \& Solusi) mengubah kesenjangan menjadi \emph{solusi}: mengelompok\-kan
gap ke dalam work package, menentukan plateau/increment, dan memilih blok bangunan
solusi. Fase F (Perencanaan Migrasi) memberi \emph{urutan dan kelayakan}:
prioritas, biaya--manfaat, risiko, dan jadwal migrasi \cite{togaf10}.
Tabel~\ref{tab:adm-phase-ef} memetakan keduanya ke kasus Bank Wanua.

\begin{table}[htbp]
\centering
\scriptsize
\caption{Pemetaan TOGAF ADM Fase E \& F ke artefak Bank Wanua.}
\label{tab:adm-phase-ef}
\begin{tabularx}{\textwidth}{@{}p{0.27\textwidth} X@{}}
\hline
\textbf{Langkah Fase E/F} & \textbf{Artefak / keluaran pada kasus Bank Wanua} \\
\hline
Analisis kesenjangan konsolidasi (E) & Gap: gateway belum ada; AML manual; BI-FAST
\& core belum terintegrasi. \\
Tentukan work package (E) & WP1 Platform API; WP2 Mesin AML; WP3 Integrasi
BI-FAST/Core; WP4 Rekonsiliasi \& Go-Live. \\
Tentukan plateau/increment (E) & Baseline $\rightarrow$ MVP Rel Transfer
$\rightarrow$ Target (SNAP BI Penuh). \\
Prioritas \& kelayakan (F) & Urutan WP, ketergantungan, biaya--manfaat, dan risiko
terhadap anggaran serta tenggat kepatuhan. \\
Rencana \& peta jalan migrasi (F) & Peta jalan berjejak: peristiwa Kickoff dan
Go-Live, deliverable per plateau. \\
\hline
\end{tabularx}
\end{table}

\subsection{Contoh Interaksi Antar-Elemen}

Gambar~\ref{fig:implementation-contoh} memakai contoh restoran: \emph{Proyek
Pemesanan Online} (work package) \emph{mewujudkan} deliverable \emph{Aplikasi
Pemesanan Online Live}, \emph{memicu} peristiwa \emph{Peluncuran}, dan
\emph{menutup} gap \emph{Belum ada pemesanan online}; deliverable itu
\emph{mewujudkan} \emph{Plateau Target Restoran} dan requirement, sedangkan
plateau \emph{mencakup} kapabilitas penyajian cepat.

\begin{figure}[htbp]
\centering
\includegraphics[width=\textwidth,height=0.4\textheight,keepaspectratio]{implementation-contoh.pdf}
\caption{Contoh interaksi antar-elemen implementasi \& migrasi pada kasus
restoran cepat saji.}
\label{fig:implementation-contoh}
\end{figure}

\subsection*{Glosarium Mini}
\begin{description}
  \item[Work package] Proyek/fase dengan awal-akhir jelas.
  \item[Deliverable] Hasil kerja terdefinisi dari work package.
  \item[Implementation event] Peristiwa implementasi (kickoff, go-live).
  \item[Plateau] Keadaan arsitektur yang stabil pada suatu periode.
  \item[Gap] Selisih antara dua plateau (hasil analisis kesenjangan).
  \item[Roadmap] Peta jalan: rangkaian plateau dan work package sepanjang waktu.
\end{description}

\section{Kerangka Berpikir}

Alur berpikir bab ini bergerak dari arsitektur target menuju rencana bertahap
(Gambar~\ref{fig:impl-flow}). Baseline dan target dibandingkan untuk menemukan
gap; gap dikelompokkan menjadi work package; work package menghasilkan deliverable;
deliverable mewujudkan plateau; plateau disusun berurutan sebagai peta jalan; dan
peta jalan menjadi bahan rekomendasi serta penutup naskah.

\begin{figure}[htbp]
\centering
\scalebox{0.86}{
\begin{tikzpicture}[
  box/.style={rectangle, rounded corners=1.5mm, draw=red!55!black,
    fill=red!8, thick, align=center, text width=30mm, minimum height=11mm},
  outnode/.style={rectangle, rounded corners=1.5mm, draw=orange!80!black,
    fill=orange!16, thick, align=center, text width=30mm, minimum height=11mm},
  link/.style={-Latex, thick, draw=red!55!black}
]
  \node[box] (gap) {Baseline vs target\\= \emph{gap}};
  \node[box, right=7mm of gap] (wp) {Work package\\per gap};
  \node[box, right=7mm of wp] (del) {Deliverable\\per work package};
  \node[box, below=8mm of del] (plat) {Plateau\\(keadaan stabil)};
  \node[outnode, left=7mm of plat] (road) {Peta jalan\\(urutan plateau)};
  \node[outnode, left=7mm of road] (rec) {Rekomendasi \&\\penutup naskah};
  \draw[link] (gap) -- (wp); \draw[link] (wp) -- (del); \draw[link] (del) -- (plat);
  \draw[link] (plat) -- (road); \draw[link] (road) -- (rec);
  \draw[link, dashed] (rec) -- (gap);
\end{tikzpicture}
}
\caption{Alur berpikir lapisan implementasi \& migrasi untuk makalah publikasi.}
\label{fig:impl-flow}
\end{figure}

\section{Konteks Indonesia}

Transformasi layanan keuangan di Indonesia jarang bisa \emph{big bang}. Sistem inti
lama, anggaran bertahap, dan tenggat kepatuhan menuntut penahapan. Konektivitas
BI-FAST dan kepatuhan format SNAP BI \cite{snapbi} kerap menjadi tonggak
(\emph{plateau}) yang menentukan urutan: kapabilitas yang berkaitan dengan
kepatuhan sering harus hidup lebih dahulu. Kewajiban penyelenggaraan sistem
elektronik---keandalan, keamanan, dan pengelolaan pusat data---membentuk kriteria
kesiapan tiap go-live \cite{pp712019}, sedangkan perlindungan data pribadi harus
sudah terpenuhi sejak plateau yang pertama kali mengolah data nasabah
\cite{uupdp2022}. Peta jalan yang realistis menempatkan tonggak kepatuhan ini
secara eksplisit, bukan sebagai renungan di akhir.

\section{Studi Kasus Utama: Peta Jalan Transformasi Bank Wanua}

Arsitektur target dari Bab~\ref{ch:strategi-bisnis}--\ref{ch:technology-architecture}
diwujudkan bertahap melalui peta jalan pada Gambar~\ref{fig:implementation-roadmap}.

\begin{figure}[htbp]
\centering
\includegraphics[width=\textwidth,height=0.5\textheight,keepaspectratio]{implementation-roadmap.pdf}
\caption{Peta jalan migrasi (ArchiMate) Bank Wanua: plateau Baseline $\rightarrow$
MVP $\rightarrow$ Target dipicu berurutan; tiap work package mewujudkan deliverable
yang mewujudkan plateau; peristiwa Kickoff dan Go-Live menandai tonggak.}
\label{fig:implementation-roadmap}
\end{figure}

\subsection{Cara Membaca Peta Jalan}

Gambar~\ref{fig:implementation-roadmap} dibaca dari kiri ke kanan sebagai garis
waktu. Tiga \textbf{plateau}---\emph{Baseline (Batch \& Manual)}, \emph{MVP Rel
Transfer}, dan \emph{Target (SNAP BI Penuh)}---dihubungkan relasi
\textbf{triggering} (memicu). Di bawahnya, empat \textbf{work package} berurutan
(WP1$\rightarrow$WP2$\rightarrow$WP3$\rightarrow$WP4); tiap work package
\textbf{mewujudkan} sebuah \textbf{deliverable}, dan deliverable itu
\textbf{mewujudkan} plateau yang sesuai (Platform API mewujudkan MVP; Mesin AML,
Adaptor, dan Rekonsiliasi mewujudkan Target). Peristiwa \emph{Kickoff Program}
memicu WP1, dan WP4 memicu \emph{Go-Live SNAP BI} yang menandai tercapainya
plateau target.

\subsection{Analisis Kesenjangan}

Peta jalan lahir dari analisis kesenjangan (Gambar~\ref{fig:implementation-gap}).
Setiap \textbf{gap} antara \emph{Baseline} dan \emph{Target}---gateway belum ada,
AML masih manual, BI-FAST \& core belum terintegrasi---\textbf{ditutup} oleh sebuah
\textbf{work package}. Membaca dari kiri ke kanan: keadaan awal (baseline), selisih
(gap), paket kerja penutup, lalu keadaan tujuan (target).

\begin{figure}[htbp]
\centering
\includegraphics[width=\textwidth,height=0.45\textheight,keepaspectratio]{implementation-gap.pdf}
\caption{Analisis kesenjangan (ArchiMate): tiap gap antara Baseline dan Target
ditutup oleh sebuah work package.}
\label{fig:implementation-gap}
\end{figure}

\subsection{Penyelarasan: Menutup Seluruh Rantai}

Karena ini lapisan terakhir, penting menunjukkan bahwa paket kerja benar-benar
menyentuh arsitektur dan mengembalikan nilai ke strategi
(Gambar~\ref{fig:implementation-stack}). Tiap \textbf{work package
\emph{membangun}} (asosiasi) komponen aplikasi yang dirancang pada
Bab~\ref{ch:application-architecture}, \textbf{mewujudkan} deliverable, dan
deliverable itu \textbf{mewujudkan} kapabilitas serta requirement dari lapisan
strategi dan motivasi (Bab~\ref{ch:togaf-lite} dan
Bab~\ref{ch:strategi-bisnis}). Dengan demikian lingkaran tertutup: motivasi menurunkan
strategi, strategi diwujudkan bisnis--aplikasi--teknologi, dan paket kerja
membangunnya kembali menjadi kapabilitas yang dijanjikan.

\begin{figure}[htbp]
\centering
\includegraphics[width=\textwidth,height=0.5\textheight,keepaspectratio]{implementation-stack.pdf}
\caption{Penyelarasan lintas-lapisan (ArchiMate): tiap work package membangun
komponen aplikasi, mewujudkan deliverable, dan deliverable mewujudkan
kapabilitas/requirement strategi--motivasi.}
\label{fig:implementation-stack}
\end{figure}

\section{Langkah Analisis dan Artefak Pendukung}

\subsection{Langkah 1: Tetapkan Baseline dan Target}
Rumuskan plateau \emph{baseline} (keadaan sekarang) dan \emph{target} (arsitektur
tujuan dari bab-bab sebelumnya).

\subsection{Langkah 2: Lakukan Analisis Kesenjangan}
Bandingkan baseline dan target, lalu nyatakan selisihnya sebagai sekumpulan gap.

\subsection{Langkah 3: Bentuk Work Package}
Kelompokkan gap menjadi work package yang koheren, masing-masing dengan deliverable
yang jelas.

\subsection{Langkah 4: Tentukan Plateau Antara}
Tetapkan satu atau lebih plateau antara (mis. MVP) yang bernilai dan dapat
dievaluasi, bukan hanya baseline dan target.

\subsection{Langkah 5: Susun Urutan dan Peristiwa}
Tentukan urutan work package (triggering), ketergantungan, dan peristiwa penting
(kickoff, go-live).

\subsection{Langkah 6: Uji Kelayakan (Fase F)}
Nilai prioritas, biaya--manfaat, dan risiko terhadap anggaran serta tenggat
kepatuhan.

\subsection{Langkah 7: Tautkan Kembali ke Strategi dan Arsitektur}
Pastikan tiap deliverable mewujudkan kapabilitas/requirement dan tiap work package
menyentuh komponen arsitektur yang sesuai.

\subsection{Langkah 8: Tulis Narasi Peta Jalan}
Ubah peta jalan menjadi paragraf yang menjelaskan urutan, tonggak, dan alasan
penahapan.

\begin{table}[htbp]
\centering
\scriptsize
\caption{Artefak pendukung Bab~9 dan fungsinya dalam makalah.}
\label{tab:artefak-bab8-impl}
\begin{tabularx}{\textwidth}{@{}p{0.26\textwidth} p{0.30\textwidth} X@{}}
\hline
\textbf{Artefak} & \textbf{Isi minimum} & \textbf{Dipakai untuk bagian makalah} \\
\hline
Katalog work package \& deliverable & Paket kerja, hasil, urutan & Rekomendasi \&
rencana. \\
Diagram peta jalan (roadmap) & Plateau, work package, peristiwa & Rencana \&
lampiran. \\
Analisis kesenjangan & Gap baseline--target \& penutupnya & Analisis \& rencana. \\
Diagram penyelarasan akhir & Work package--deliverable--kapabilitas & Sintesis \&
keterlacakan end-to-end. \\
\hline
\end{tabularx}
\end{table}

\section{Integrasi ke Makalah Publikasi}

Lapisan ini mengisi bagian \textbf{Rekomendasi/Rencana} dan menutup naskah. Peta
jalan menjadi bukti bahwa rekomendasi arsitektur dapat dieksekusi secara realistis
dan patuh. Contoh paragraf:

\begin{quote}
Transformasi diusulkan dalam tiga plateau. Plateau MVP memasang platform API dan
gateway sehingga transfer real-time dasar berjalan; plateau target menambahkan AML
otomatis, adaptor BI-FAST \& core, serta rekonsiliasi hingga patuh SNAP BI penuh.
Penahapan ini menautkan tiap paket kerja dengan kapabilitas strategis yang
diwujudkannya, sehingga anggaran terbatas dialokasikan pada tonggak yang paling
bernilai dan paling terikat kepatuhan lebih dahulu.
\end{quote}

\section{Diskusi Kritis, Trade-off, dan Perangkap Umum}

\begin{table}[htbp]
\centering
\scriptsize
\caption{Trade-off dalam perencanaan implementasi \& migrasi.}
\label{tab:tradeoff-bab8-impl}
\begin{tabularx}{\textwidth}{@{}p{0.24\textwidth} X X@{}}
\hline
\textbf{Keputusan} & \textbf{Jika terlalu ringan} & \textbf{Jika terlalu berat} \\
\hline
Jumlah plateau & Satu lompatan besar, berisiko & Terlalu banyak tonggak,
lambat \& mahal. \\
Ukuran work package & Terlalu besar, sulit dikelola & Terlalu kecil, overhead
koordinasi. \\
Urutan \& kepatuhan & Kepatuhan tertinggal di akhir & Semua sekaligus, anggaran
jebol. \\
Sistem inti lama & Diabaikan, integrasi rapuh & Diganti sekaligus, risiko
tinggi. \\
\hline
\end{tabularx}
\end{table}

Perangkap umum: (1) menggambar peta jalan tanpa analisis kesenjangan yang jujur;
(2) plateau yang tidak bernilai (sekadar milestone teknis); (3) melupakan
peristiwa kepatuhan (go-live SNAP BI) sebagai penentu urutan; dan (4) work package
yang tidak tertaut ke kapabilitas mana pun.

\section{Latihan Terapan Individual}

\subsection*{Latihan~9.1: Baseline, Target, dan Gap}
Untuk organisasi studi kasus Anda, rumuskan plateau baseline dan target, lalu
daftarkan gap di antaranya.

\subsection*{Latihan~9.2: Work Package dan Deliverable}
Kelompokkan gap menjadi work package, tentukan deliverable tiap paket, dan tautkan
ke kapabilitas/requirement yang diwujudkan.

\subsection*{Latihan~9.3: Peta Jalan Berplateau}
Susun peta jalan dengan minimal satu plateau antara, urutan work package, dan
peristiwa penting (kickoff/go-live).

\subsection*{Latihan~9.4: Paragraf Rencana Migrasi}
Tulis satu paragraf yang menjelaskan urutan, tonggak, dan alasan penahapan peta
jalan Anda.

\begin{table}[htbp]
\centering
\scriptsize
\caption{Rubrik latihan Bab~9.}
\label{tab:rubrik-bab8-impl}
\begin{tabularx}{\textwidth}{@{}p{0.25\textwidth} X X@{}}
\hline
\textbf{Komponen} & \textbf{Indikator baik} & \textbf{Perlu diperbaiki bila} \\
\hline
Baseline--target--gap & Gap konkret \& berbasis selisih nyata & Gap kabur atau
sekadar daftar keinginan. \\
Work package \& deliverable & Paket koheren, deliverable jelas \& tertaut & Paket
tanpa hasil atau tanpa tautan kapabilitas. \\
Peta jalan & Plateau bernilai, urutan \& tonggak jelas & Satu lompatan atau
milestone teknis tanpa nilai. \\
Narasi rencana & Menjelaskan urutan \& alasan penahapan & Deskriptif tanpa
argumen kelayakan. \\
\hline
\end{tabularx}
\end{table}

\section{Ringkasan, Refleksi, dan Penutup Buku}

\subsection*{Poin Kunci}
\begin{itemize}
  \item Lapisan implementasi \& migrasi memuat work package, deliverable,
        implementation event, plateau, dan gap.
  \item \textbf{Work package mewujudkan deliverable}; \textbf{deliverable
        mewujudkan plateau}; \textbf{plateau mencakup} kapabilitas/requirement;
        \textbf{gap} menandai selisih baseline--target.
  \item Peta jalan menyusun plateau secara berurutan (triggering) dengan peristiwa
        penting sebagai tonggak.
  \item Fase E mengubah gap menjadi work package \& plateau; Fase F memberi urutan,
        biaya, dan risiko.
  \item Penahapan yang realistis menautkan tiap paket kerja ke kapabilitas dan
        menempatkan tonggak kepatuhan lebih dahulu.
\end{itemize}

\subsection*{Uji Mandiri}
\begin{enumerate}
  \item Jelaskan perbedaan plateau, gap, work package, dan deliverable.
  \item Mengapa migrasi dinyatakan sebagai rangkaian plateau, bukan satu lompatan
        dari baseline ke target?
  \item Ambil satu gap pada kasus Anda dan sebutkan work package, deliverable, dan
        plateau yang menutupnya.
\end{enumerate}

\subsection*{Penutup: Menutup Rantai Tujuh Lapisan}
Buku ini menelusuri arsitektur enterprise dari \textbf{motivasi}
(Bab~\ref{ch:strategi-bisnis})---mengapa---turun ke \textbf{strategi}
(Bab~\ref{ch:togaf-lite}), \textbf{bisnis}
(Bab~\ref{ch:business-architecture}), \textbf{data}
(Bab~\ref{ch:data-architecture}), \textbf{aplikasi}
(Bab~\ref{ch:application-architecture}), dan \textbf{teknologi}
(Bab~\ref{ch:technology-architecture})---apa dan bagaimana secara struktural---lalu
ditutup pada bab ini dengan \textbf{implementasi \& migrasi}---bagaimana
mencapainya sepanjang waktu. Setiap lapisan menurun dan mewujudkan lapisan di
atasnya, dan pada akhirnya paket kerja membangunnya kembali menjadi kapabilitas
yang dijanjikan motivasi. Inilah kekuatan ArchiMate yang dipadu TOGAF ADM: satu
model yang dapat ditelusuri dari ``mengapa'' sampai ``kapan dan dengan apa''.
Mahasiswa yang menuntaskan ketujuh lapisan pada satu studi kasus telah menulis
sebuah arsitektur enterprise yang utuh, koheren, dan siap dipertahankan secara
akademik.

\section{Bacaan Lanjutan dan Lampiran}

\subsection*{Bacaan Inti}
\begin{itemize}
  \item The Open Group untuk Fase E \& F ADM dan artefak implementasi \& migrasi
        \cite{togaf10}.
  \item Ross, Weill, \& Robertson untuk fondasi eksekusi dan urutan investasi
        kapabilitas \cite{ross2006enterprise}.
\end{itemize}

\subsection*{Bacaan untuk Memperdalam}
\begin{itemize}
  \item Newman untuk pengiriman inkremental dan penggelaran bertahap
        \cite{newman2021building}.
  \item Untuk konteks Indonesia: standar API pembayaran \cite{snapbi},
        penyelenggaraan sistem elektronik \cite{pp712019}, dan perlindungan data
        pribadi \cite{uupdp2022}.
\end{itemize}

\subsection*{Lampiran Kerja Bab~9}
\begin{description}
  \item[Katalog work package \& deliverable] Paket kerja, hasil, urutan.
  \item[Diagram peta jalan] Plateau, work package, peristiwa.
  \item[Analisis kesenjangan] Gap baseline--target dan penutupnya.
  \item[Diagram penyelarasan akhir] Work package--deliverable--kapabilitas.
\end{description}
\setlength{\emergencystretch}{0pt}
